% JNE-specific replacement of original BSPC editorial comparison.
\subsection{Relation to established neural-engineering validation studies}
\label{sec:p1-jne-peer}
Recent accepted research in Journal of Neural Engineering shows
that a neural decoding methods paper need not report a new SOTA
network, but does need a generalisable, experimentally resolved
measurement or methodological question.
Eder et al.\ \cite{eder2024jnebenchmark} compared Riemannian
decoding against EEGNet and other convolutional models
across motor imagery, P300 and SSVEP paradigms with multiple
validation protocols and reproducible open-source pipelines.
Seibel et al.\ \cite{seibel2026jnedata}
measured how CNN decoding performance responds to subject
count and recorded EEG duration across three tasks.
Their conclusions are operational: they directly inform
model family or data acquisition decisions.
Our scope is narrower---binary motor-imagery
tangent-space linear discriminants and the
selection \emph{criterion}, not a deep learning
architecture leaderboard nor a data-acquisition
learning curve.

\begin{table}[tb]
\centering\small
\caption{Comparison with representative published
JNE methods studies, preserving non-comparable
endpoints. This is a design/evidence comparison,
not an accuracy leaderboard.}
\label{tab:p1-jne-peer}
\setlength{\tabcolsep}{4pt}
\begin{tabular}{p{.24\linewidth}p{.31\linewidth}p{.33\linewidth}}
\toprule
Study & Experimental scope & Actionable conclusion\\
\midrule
Eder et al.\ \cite{eder2024jnebenchmark}
& MI/P300/SSVEP; CNN and Riemannian; multiple split budgets
& Comparative model-family tradeoffs\\
Seibel et al.\ \cite{seibel2026jnedata}
& Three EEG tasks and CNN structures; varied training sample supply
& Estimate marginal data-acquisition value\\
This study
& Four original frozen AUC experiment strata; three
  exposed source-only EEG score-anatomy collections
& Separate genuinely predictive from class-centred
  source scoring and disclose invariant penalties;
  no new optimal transfer selector established\\
\bottomrule
\end{tabular}
\end{table}

We aim to supply an equally \emph{verifiable} conclusion
within our constrained decoder family: a historical
validation covariance score can be misidentified as
source-trained class-conditional predictive likelihood,
and even a correctly computed Gaussian score can be
dominated by source validation shifts invisible to
a fixed within-person rank metric.
The scientific value lies in the distinction
between these estimands and the downstream
interpretation of historical source choices.
Unlike the two cited JNE examples, we have
\emph{not} demonstrated that a replacement
scoring policy gives consistently better
future-participant decoding. That limitation
is an explicit boundary of the manuscript.

